\documentclass[11pt]{article}
\newcommand{\paso}{5}
\usepackage[T1]{fontenc}
\usepackage[utf8]{inputenc}
\usepackage[spanish,es-nodecimaldot,es-noquoting,es-lcroman]{babel}
\spanishplainpercent
\usepackage{lmodern}
\usepackage[a4paper,margin=21mm,headheight=15pt]{geometry}
\usepackage{amsmath,amssymb,bm}
\usepackage{booktabs,array}
\usepackage{tikz}
\usetikzlibrary{arrows.meta,positioning,calc}
\usepackage[most]{tcolorbox}
\usepackage{microtype}
\usepackage{fancyhdr}
\usepackage{hyperref}
\hypersetup{colorlinks=true,linkcolor=blue!45!black,urlcolor=blue!45!black}
\setlength{\parindent}{0pt}
\setlength{\parskip}{5pt}
\setlength{\emergencystretch}{2em}
\pagestyle{fancy}
\fancyhf{}
\fancyhead[L]{\small BEM del cilindro: paso \paso}
\fancyhead[R]{\small Guía matemática independiente}
\fancyfoot[C]{\thepage}
\newcommand{\vct}[1]{\bm{#1}}
\newcommand{\dd}{\,\mathrm{d}}
\newcommand{\ii}{\mathrm{i}}
\newcommand{\norm}[1]{\left\lVert#1\right\rVert}
\newcommand{\VP}{\operatorname{VP}}
\newcommand{\tr}{\operatorname{tr}}
\newtcolorbox{idea}[1]{breakable,colback=blue!3,colframe=blue!50!black,
  title={#1},fonttitle=\bfseries,boxrule=0.7pt}
\newtcolorbox{alerta}[1]{breakable,colback=orange!4,colframe=orange!65!black,
  title={#1},fonttitle=\bfseries,boxrule=0.7pt}
\newtcolorbox{resultado}[1]{breakable,colback=teal!3,colframe=teal!55!black,
  title={#1},fonttitle=\bfseries,boxrule=0.7pt}
\newcommand{\inicio}[2]{%
 {\small\color{blue!55!black}\bfseries SERIE: LOS SIETE PASOS MATEMÁTICOS DEL BEM}\par
 {\LARGE\bfseries Paso \paso. #1\par}\vspace{3mm}
 \begin{idea}{¿Qué busca este paso?}#2\end{idea}}
\newcommand{\mapa}[3]{%
 \begin{center}
 \begin{tikzpicture}[node distance=4mm,
   caja/.style={draw=blue!50!black,fill=blue!3,rounded corners=2pt,
     text width=13cm,align=center,inner sep=6pt},
   >={Stealth[length=2mm]}]
 \node[caja] (a){#1};
 \node[caja,below=of a] (b){#2};
 \node[caja,below=of b] (c){#3};
 \draw[->,thick,blue!60!black] (a)--(b);
 \draw[->,thick,blue!60!black] (b)--(c);
 \end{tikzpicture}
 \end{center}}
\newcommand{\base}{%
 \smallskip{\footnotesize\color{black!65}
 Modelo de referencia: cilindro sólido, material homogéneo e isótropo,
 deformación plana, pequeñas deformaciones, sin fuerzas de volumen;
 amplitudes complejas con \(e^{+\ii\omega t}\), \(\omega>0\).}\par}
\newcommand{\fuentes}{%
 \par\medskip\begingroup\footnotesize
 \textbf{Referencia y continuación.}
 Este documento desarrolla una etapa del esquema
 \texttt{esquema\_ecuaciones\_bem.tex}. Las ecuaciones y los ejemplos se
 contrastaron con \texttt{bem\_cilindro.py}, ambos en la carpeta superior.
 La guía general y sus deducciones ampliadas se encuentran en
 \texttt{../guia\_bem/guia\_bem.pdf}. No se modifica la implementación.
 \par Para recompilar este documento, mantener \texttt{estilo.tex} en
 la misma carpeta y ejecutar \texttt{pdflatex} sobre su archivo
 \texttt{.tex}, desde esa carpeta.\par\endgroup}

\begin{document}
\inicio{Discretizar e integrar}{Convertir una frontera continua y sus campos
en elementos, valores nodales y sumas de cuadratura. Este paso decide
cómo se representa la geometría, cómo varían los campos entre nodos y
dónde se evalúan los núcleos del paso 4. Todavía no resuelve el sistema.}
\base

\mapa{Curva continua \(\Gamma\), desplazamiento \(\vct u(\vct x)\), tracción \(\vct t(\vct x)\)}
{Elementos de tres nodos + funciones de forma cuadráticas en \(\eta\in[-1,1]\)}
{Puntos \(\vct x_q\), normales \(\vct n_q\), jacobianos \(J_q\) y pesos \(W_{qa}\)}

\section*{1. Qué significa discretizar el contorno}
Se divide la circunferencia en \(n_e\) elementos. Cada elemento tiene un nodo
en cada extremo y otro intermedio. Los elementos vecinos comparten sus
extremos, y el último se conecta con el primero.
Por eso hay \(N=2n_e\) nodos globales, no \(3n_e\).

\begin{center}
\begin{tikzpicture}[>={Stealth[length=2mm]}]
 \draw[thick,blue!60!black] (-6,0)--(-2,0);
 \foreach \x/\lab in {-6/-1,-4/0,-2/1}{
   \fill (\x,0) circle(2pt);
   \node[below] at (\x,0){$\eta=\lab$};}
 \node at (-4,1){Elemento de referencia};
 \draw[->,thick] (-1.6,0)--(-0.4,0);
 \draw[thick,teal!60!black] (0,-0.2) .. controls (1.3,1) and (2.7,1) .. (4,-0.2);
 \fill (0,-0.2) circle(2pt) node[below]{$\vct x_1$};
 \fill (2,0.7) circle(2pt) node[above]{$\vct x_2$};
 \fill (4,-0.2) circle(2pt) node[below]{$\vct x_3$};
 \node at (2,1.5){Elemento físico};
\end{tikzpicture}
\end{center}
El dibujo muestra el concepto de mapeo. En el modelo se usan polinomios
cuadráticos para interpolar las coordenadas del círculo.

\section*{2. Coordenadas, conectividad y componentes}
Con índices globales \(g=0,\ldots,N-1\):
\[
 \theta_g=\frac{2\pi g}{N},\qquad
 \vct x_g=R(\cos\theta_g,\sin\theta_g).
\]
Para el elemento \(e=0,\ldots,n_e-1\), los tres índices globales son
\[
 g(e,1)=2e,\qquad g(e,2)=2e+1,\qquad
 g(e,3)=(2e+2)\bmod N.
\]
El índice local \(a=1,2,3\) no es un índice global.
Cada nodo tiene dos componentes: con ocho elementos hay 16 nodos y
32 incógnitas escalares de desplazamiento.

\newpage
\section*{3. Deducir las funciones de forma}
Se buscan polinomios \(P_a\) de grado dos que valgan uno en su propio nodo y
cero en los otros. Para los nodos locales \(\eta_1=-1,\eta_2=0,\eta_3=1\),
la fórmula de Lagrange es
\[
 P_a(\eta)=\prod_{b\ne a}\frac{\eta-\eta_b}{\eta_a-\eta_b}.
\]
Al sustituir los tres valores:
\[
 \boxed{P_1=\tfrac12\eta(\eta-1),\qquad
 P_2=1-\eta^2,\qquad P_3=\tfrac12\eta(\eta+1).}
\]
En el esquema general se llamaban \(N_a\); aquí usamos \(P_a\) para no
confundirlas con el número \(N\) de nodos globales.

\begin{center}
\begin{tabular}{c|ccc}
\toprule
 & \(\eta=-1\)&\(\eta=0\)&\(\eta=1\)\\
\midrule
\(P_1\)&1&0&0\\
\(P_2\)&0&1&0\\
\(P_3\)&0&0&1\\
\bottomrule
\end{tabular}
\end{center}
\[
 P_1+P_2+P_3=1,\qquad
 P_1'=\eta-\tfrac12,\quad P_2'=-2\eta,\quad
 P_3'=\eta+\tfrac12.
\]
La suma uno es la partición de la unidad: permite representar exactamente
un campo constante y será esencial para la traslación rígida del paso 6.

\section*{4. Interpolar geometría y campos con la misma base}
\[
\begin{aligned}
 \vct x_e(\eta)&=\sum_{a=1}^3P_a(\eta)\vct x_{g(e,a)},\\
 \vct u_e(\eta)&\simeq\sum_{a=1}^3P_a(\eta)\vct u_{g(e,a)},\\
 \vct t_e(\eta)&\simeq\sum_{a=1}^3P_a(\eta)\vct t_{g(e,a)}.
\end{aligned}
\]
Se llama interpolación \emph{isoparamétrica} porque geometría y campos usan
las mismas funciones. Los valores nodales no dependen de \(\eta\);
por eso podrán salir de las integrales al ensamblar.

\begin{alerta}{Interpolar no significa que todo sea un promedio positivo}
Por ejemplo, en \(\eta=1/2\):
\[
 (P_1,P_2,P_3)=(-1/8,\,3/4,\,3/8).
\]
Una función de forma puede ser negativa y aun así interpolar correctamente.
\end{alerta}

\newpage
\section*{5. Tangente, jacobiano y normal}
La derivada de la curva interpolada es
\[
 \vct x'_e(\eta)=\sum_aP_a'(\eta)\vct x_{g(e,a)},\qquad
 J=\norm{\vct x'_e},\qquad
 \dd\Gamma=J\,\dd\eta.
\]
Como \(\eta\) es adimensional, \(J\) tiene unidades de longitud.
La tangente unitaria y la normal exterior, para recorrido antihorario, son
\[
 \vct\tau=\frac{(x',y')}{J},\qquad
 \boxed{\vct n=\frac{(y',-x')}{J}.}
\]
La normal es la tangente rotada \(90^\circ\) hacia la derecha.
Al invertir la conectividad, la misma fórmula da la normal opuesta.

\section*{6. Ejemplo geométrico del primer elemento}
Para \(R=6\) m y ocho elementos:
\[
\begin{aligned}
 \vct x_1&=(6,0),\\
 \vct x_2&\simeq(5.543277,\,2.296101),\\
 \vct x_3&\simeq(4.242641,\,4.242641).
\end{aligned}
\]
En \(\eta=0\), \(P_2=1\), por lo que la posición es exactamente
\(\vct x_2\). Además,
\[
 \vct x'(0)=\tfrac12(\vct x_3-\vct x_1),\quad
 J(0)\simeq2.296101\ \mathrm m,\quad
 \vct n(0)\simeq(0.923880,\,0.382683).
\]
En \(\eta=1/2\), la interpolación da
\[
 \vct x(1/2)=-\tfrac18\vct x_1+\tfrac34\vct x_2+\tfrac38\vct x_3
 \simeq(4.998448,\,3.313066)\ \mathrm m.
\]

\begin{alerta}{Los nodos están en el círculo, la curva entre ellos es aproximada}
La interpolación polinómica cuadrática no representa exactamente un arco
circular. En el extremo inicial de este elemento, la normal geométrica es
aproximadamente \((0.999901,0.014069)\), mientras que la normal circular
nodal es \((1,0)\). La diferencia disminuye al refinar la malla.
\end{alerta}

\newpage
\section*{7. Convertir una integral en una suma}
Para una función regular:
\[
 \int_{\Gamma_e}f(\vct x)\dd\Gamma
 =\int_{-1}^{1}f(\vct x_e(\eta))J(\eta)\dd\eta
 \simeq\sum_{q=1}^{n_q}f(\vct x_q)J_qw_q.
\]
Los puntos \(\eta_q\) y pesos \(w_q\) de Gauss-Legendre no son uniformes.
Para el aporte asociado al nodo local \(a\), conviene agrupar
\[
 \boxed{W_{qa}=P_a(\eta_q)J_qw_q.}
\]
De ese modo,
\[
 \int_{\Gamma_e}P_aU\,\dd\Gamma\simeq\sum_qW_{qa}U_q,\qquad
 \int_{\Gamma_e}P_a(T-T^s)\,\dd\Gamma\simeq\sum_qW_{qa}(T_q-T_q^s).
\]
Los pesos no contienen el núcleo: se calculan a partir de geometría,
funciones de forma y cuadratura.

\section*{8. Si la fuente pertenece al elemento}
Si el nodo de colocación es uno de sus tres nodos, la fuente está en
\(\eta_s=-1,0\) o \(1\). Se divide el intervalo en
\([-1,\eta_s]\) y \([\eta_s,1]\), descartando tramos de longitud cero.
En un tramo \([a,b]\), usando \(t\in(0,1)\):
\[
\begin{array}{ll}
\text{Fuente a la izquierda:}&
\eta=a+(b-a)t^2,\quad \dd\eta=2(b-a)t\,\dd t,\\[1mm]
\text{Fuente a la derecha:}&
\eta=a+(b-a)[1-(1-t)^2],\\
&\dd\eta=2(b-a)(1-t)\,\dd t.
\end{array}
\]
Así se concentran puntos cerca de la fuente y el nuevo peso incluye la
derivada del mapa.

\begin{resultado}{Por qué ayuda para la singularidad logarítmica}
Si \(\eta-a\propto t^2\), un factor
\(\log|\eta-a|\dd\eta\) se transforma en un comportamiento proporcional a
\(t\log t\,\dd t\), cuyo límite es cero. El integrando resulta más favorable,
pero no se vuelve necesariamente un polinomio ni una función analítica.
\end{resultado}

\newpage
\section*{9. Cuántos puntos se usan en este caso}
\begin{center}
\begin{tabular}{p{0.52\linewidth}p{0.34\linewidth}}
\toprule
Posición de la fuente respecto al elemento & Regla utilizada\\
\midrule
Fuera del elemento & 16 puntos de Gauss ordinarios.\\
En un extremo & Un subintervalo transformado, 16 puntos.\\
En el nodo intermedio & Dos subintervalos transformados, 32 puntos en total.\\
\bottomrule
\end{tabular}
\end{center}
Los puntos no coinciden con la singularidad. Pero esto no basta para
definir el valor principal del término estático \(T^s\): se necesita
la construcción diagonal del paso 6.

\section*{10. Tres normales y campos que no deben confundirse}
\begin{enumerate}
 \item \textbf{Normal de cuadratura:} deriva de la curva interpolada y entra
 en \(T\) y \(T^s\).
 \item \textbf{Normal circular nodal:} \(\vct n_g=\vct x_g/R\);
 se usa para fijar \(\vct t_g=-p\vct n_g\) y proyectar \(\vct u_g\).
 \item \textbf{Tracción interpolada:}
 \(\vct t_e(\eta)=\sum_aP_a(-p\vct x_a/R)\).
 En general no coincide exactamente con \(-p\vct n_e(\eta)\).
\end{enumerate}
No es una contradicción en las definiciones, sino una aproximación
geométrica y de interpolación del esquema existente.

\section*{11. Comprobaciones antes de ensamblar}
\begin{itemize}
 \item \(P_a(\eta_b)=\delta_{ab}\), \(\sum_aP_a=1\) y \(\sum_aP_a'=0\).
 \item \(J>0\), \(\norm{\vct n}=1\), \(\vct n\cdot\vct\tau=0\);
 comprobar también que la normal apunte hacia afuera.
 \item La conectividad debe cerrar el contorno sin duplicar indebidamente
 el último nodo.
 \item Los 16 puntos integran exactamente polinomios de grado hasta 31,
 pero los núcleos no son polinomios: no hay exactitud automática.
 \item Aumentar puntos no sustituye refinar elementos. Son errores
 distintos; además la cancelación en los núcleos limita la precisión.
\end{itemize}
\begin{resultado}{Qué queda preparado}
Para cada par fuente--elemento se dispone de
\(\{\vct x_q,\vct n_q,W_{qa}\}\).
Con ellos, el paso 6 transforma las integrales de influencia en bloques
numéricos y los suma en el lugar global correspondiente.
\end{resultado}
\fuentes
\end{document}
